\chapter{Istniejące algorytmy planowania czasu procesora w nowoczesnych systemach}
\label{ch:istniejace}

\section{System Linux i architektura EEVDF}
W nowoczesnych wersjach jądra systemu Linux, domyślnym algorytmem planowania (ang. \textit{scheduler}) dla procesów standardowych jest algorytm Earliest Eligible Virtual Deadline First (EEVDF). Zastąpił on wieloletni standard znany jako Completely Fair Scheduler (CFS), rozwiązując fundamentalny problem balansu między sprawiedliwością przydziału czasu a opóźnieniami w obsłudze zadań interaktywnych \cite{linuxdocEEVDF, lkmlSchedV612}. 

Działanie algorytmu EEVDF opiera się na ścisłych podstawach matematycznych \cite{stoica1995eevdf} i operuje na kilku kluczowych wartościach przypisywanych każdemu zadaniu:
\begin{itemize}
	\item \textbf{Żądany kwant czasu} (\textit{requested time slice / slice})
	\item \textbf{Wirtualny czas wykonania} (\textit{virtual runtime / vruntime})
	\item \textbf{Opóźnienie} (\textit{lag})
	\item \textbf{Priorytet} (\textit{priority / weight})
\end{itemize}

Kiedy nowe lub powracające ze stanu uśpienia zadanie trafia do kolejki procesów gotowych, otrzymuje bazowy kwant czasu (\textit{slice}). Wartość ta oparta jest na zmiennych konfiguracyjnych systemu (np. \texttt{sysctl\_sched\_base\_slice}) i jest dynamicznie skalowana na podstawie wagi procesu (wynikającej z jego priorytetu \textit{nice}) \cite{linuxFairC, lkmlRunToParity}. Następnie, gdy zadanie zaczyna być wykonywane przez procesor, jego wirtualny czas wykonania (\texttt{vruntime}) stale rośnie. Wartość ta odzwierciedla rzeczywisty czas spędzony na procesorze, znormalizowany przez wagę zadania. Zliczanie tego czasu opiera się na precyzyjnych licznikach sprzętowych procesora (np. w architekturze amd64), jednak sama aktualizacja wirtualnego czasu i struktur danych jest w całości realizowana programowo przez jądro systemu operacyjnego, zazwyczaj podczas obsługi przerwań zegarowych \cite{linuxSchedH, linuxFairC}.

Na bazie początkowo przypisanego kwantu czasu oraz \texttt{vruntime}, system wylicza wirtualny termin zakończenia (\textit{virtual deadline}). Reprezentuje on moment, do którego proces powinien w idealnych warunkach skonsumować swój przydzielony czas. Algorytm na bieżąco oblicza również opóźnienie (\textit{lag}), będące różnicą pomiędzy czasem, jaki zadanie powinno było otrzymać w idealnym, proporcjonalnym modelu podziału czasu (tzw. \textit{fluid-flow model}), a czasem, który faktycznie wykorzystało. Jeżeli wartość \textit{lag} jest mniejsza od zera, oznacza to, że proces otrzymał więcej czasu procesora niż pozostałe procesy rywalizujące o czas procesora. W takim scenariuszu proces traci kwalifikowalność (\textit{ineligible}) i jest kategorycznie wykluczony z możliwości przydziału sprzętowego do momentu, aż ogólny upływ czasu w systemie sprawi, że jego \textit{lag} ponownie wzrośnie do wartości nieujemnej \cite{linuxdocEEVDF, stoica1995eevdf}.

\subsection{Wybór procesu i struktury danych}
Działanie każdego planisty w nowoczesnych systemach operacyjnych opiera się na sprzętowym mechanizmie przerwań czasowych (ang. \textit{timer interrupts}). Ponieważ procesor nie może jednocześnie wykonywać własnego kodu schedulera oraz kodu aplikacji użytkownika, wbudowany kontroler przerwań z dużą częstotliwością wstrzymuje wykonywane zadanie, przekazując sterowanie do jądra systemu \cite{wikiScheduling}. Daje to algorytmowi możliwość przeliczenia wskaźników i, w razie potrzeby, wywłaszczenia procesu.

Scheduler EEVDF przechowuje informacje o aktywnych procesach w zrównoważonym drzewie czerwono-czarnym (ang. \textit{Augmented Red-Black Tree}). W przeciwieństwie do naiwnego podejścia, drzewo to nadal sortuje węzły na podstawie wirtualnego czasu wykonania (\texttt{vruntime}), co układa zadania od lewej do prawej. Kluczową innowacją jest jednak rozszerzenie (\textit{augmentation}) każdego węzła o dodatkowe pole \texttt{min\_deadline}, które przechowuje najkrótszy termin zakończenia z danego węzła i wszystkich jego potomków. Dzięki temu struktura działa jednocześnie jako drzewo wyszukiwania binarnego oraz jako kopiec (min-heap) dla wartości \textit{deadline}. Architektura ta pozwala systemowi w optymalnym czasie $O(\log n)$ odrzucić całe poddrzewa i natychmiast odnaleźć zadanie o najwcześniejszym \textit{virtual deadline}, które jednocześnie spełnia warunek kwalifikowalności ($lag \ge 0$) i to ono zostaje wybrane do działania na procesorze \cite{kernelInternalsEEVDF, linuxFairC}.

\subsection{Mechanizm Run to Parity}
Algorytm EEVDF w swojej czystej, teoretycznej postaci zakłada wyidealizowany model maszyny obliczeniowej. W rzeczywistości fizyczne procesory opierają swoją wydajność na hierarchii pamięci podręcznej (Cache L1/L2/L3), która przechowuje dane powiązane z aktualnie wykonywanym procesem. Niezmodyfikowany mechanizm doboru zadań oparty wyłącznie na \textit{deadline'ach} prowadziłby do zjawiska tzw. \textit{over-scheduling} – zbyt częstych zmian aktualnie działającego zadania. Skutkowałoby to znacznymi stratami cykli procesora na ciągłą wymianę linii pamięci podręcznej (\textit{cache miss}) oraz narzutem związanym z przełączaniem kontekstu.

Aby rozwiązać ten problem, w systemie Linux wprowadzono mechanizm \textit{Run to Parity} \cite{lkmlRunToParity}. Zgodnie z tą regułą, proces, który już przejął kontrolę nad jednostką obliczeniową, z reguły nie zostaje wywłaszczony przez inne, standardowe zadania posiadające wcześniejszy \textit{virtual deadline}, dopóki jego własne opóźnienie (\textit{lag}) nie spadnie do zera. Dopiero gdy proces "wyrówna" swój dług wobec systemu (osiągając parytet), zdejmuje się z niego ochronę, pozwalając na przejęcie procesora przez kolejne kwalifikujące się zadanie. Mechanizm ten gwarantuje dłuższe, nieprzerwane odcinki pracy pojedynczego wątku, co maksymalizuje wykorzystanie pamięci podręcznej i zwiększa ogólną przepustowość sprzętu (throughput), nie naruszając jednocześnie długoterminowych matematycznych gwarancji sprawiedliwości algorytmu \cite{lkmlSchedV612}.






\section{Systemy z grupy Windows}
W nowoczesnych systemach z rodziny Windows za przydział czasu procesora odpowiada planista (ang. \textit{scheduler}) zaprojektowany pierwotnie na potrzeby architektury Microsoft Windows NT. W przeciwieństwie do algorytmów opartych na sprawiedliwym podziale czasu, rozwiązanie firmy Microsoft bazuje na wywłaszczeniowym planowaniu wielozadaniowym, w którym absolutnym fundamentem jest rygorystyczny priorytet \cite{loyolaProcThreadScheduling}. Co istotne, jednostką planowania w systemie Windows nie jest proces, lecz pojedynczy wątek. Proces pełni jedynie rolę kontenera na zasoby (takie jak pamięć czy kontekst zabezpieczeń), podczas gdy o dostęp do jednostki obliczeniowej rywalizują bezpośrednio wątki \cite{msLearnProcessesThreads}.

Architektura planisty definiuje 32 odrębne poziomy priorytetów (od 0 do 31), które podzielone są na trzy główne kategorie \cite{solomon2000threadScheduling}:
\begin{itemize}
	\item \textbf{Poziom 0 (Zero-Page Thread):} Zarezerwowany wyłącznie dla systemowego wątku bezczynności (\textit{idle thread}), który uruchamia się tylko przy braku innych zadań. Zajmuje się on m.in. asynchronicznym zerowaniem wolnych stron pamięci fizycznej.
	\item \textbf{Poziomy 1--15 (Priorytety dynamiczne):} Przydzielane standardowym wątkom aplikacji użytkownika oraz niekrytycznym wątkom jądra. Podlegają one ciągłym, dynamicznym algorytmom modyfikacji (wzmacnianiu i degradacji).
	\item \textbf{Poziomy 16--31 (Priorytety czasu rzeczywistego):} Zarezerwowane dla absolutnie krytycznych operacji (np. przetwarzanie multimediów przy minimalnych opóźnieniach, sterowniki sprzętowe). Wątki na tych poziomach mają sztywny priorytet i nie podlegają dynamicznym zmianom \cite{msLearnPriorityBoosts}.
\end{itemize}

Przy tworzeniu wątku w paśmie dynamicznym przydzielana jest mu tzw. wartość bazowa priorytetu (\textit{base priority}), obliczana na podstawie klasy priorytetu procesu oraz priorytetu wątku wewnątrz tej klasy. Aby zapewnić wysoką responsywność interaktywną i zapobiegać zbyt długiemu oczekiwaniu zadań (\textit{starvation}), planista stosuje mechanizm dynamicznego wzmacniania priorytetu (\textit{priority boosting}). W systemie istnieje sześć głównych rodzajów zdarzeń, które wyzwalają tymczasowe podbicie priorytetu wątku \cite{bsodTutorialsQuantumBoost}:
\begin{enumerate}
	\item Zakończenie asynchronicznej operacji wejścia/wyjścia (I/O completion).
	\item Wybudzenie wątku przez obiekty synchronizacji (np. semafory, zdarzenia).
	\item Zakończenie oczekiwania na zasoby warstwy wykonawczej (Executive resources).
	\item Działanie w ramach procesu pierwszoplanowego (Foreground process).
	\item Wybudzenie w celu obsługi zdarzeń interfejsu graficznego (GUI windowing activity).
	\item Interwencja usługi \textit{Balance Set Manager}, chroniąca przed głodzeniem poprzez wywindowanie priorytetu zablokowanego wątku do wartości 15 \cite{duffyBalanceSetManager}.
\end{enumerate}

Po otrzymaniu wzmocnienia, kwant czasu (ang. \textit{quantum}) wątku ulega wyzerowaniu. Kiedy wątek w pełni wykorzysta nowo przydzielony kwant, system stosuje mechanizm degradacji (\textit{priority decay}) – jego priorytet jest zmniejszany o 1. Proces ten powtarza się przy każdym kolejnym wyczerpaniu kwantu, aż wartość spadnie z powrotem do nominalnego priorytetu bazowego \cite{bsodTutorialsQuantumBoost}.

\subsection{Wybór procesu i struktury danych}
Kernel utrzymuje niezależne struktury danych dla każdego procesora logicznego, zgrupowane wewnątrz bloku \textit{Kernel Processor Control Block} (KPRCB) \cite{reactosThrdschdC}. Serce tej struktury stanowią kolejki gotowości (\textit{Dispatcher Ready Queues}), zaimplementowane jako tablica 32 dwukierunkowych list powiązanych – po jednej dla każdego poziomu priorytetu. Nowe lub wybudzone wątki dodawane są na koniec listy odpowiadającej ich priorytetowi.

W celu optymalizacji wyboru następnego zadania, dyspozytor unika czasochłonnego iterowania po listach. Utrzymuje on pojedynczą, 32-bitową liczbę, zwaną \textit{Ready Summary}. Każdy bit tej zmiennej powiązany jest z jedną kolejką. Jeśli dana kolejka posiada przynajmniej jeden wątek gotowy do natychmiastowego uruchomienia, odpowiadający jej bit ustawiany jest na wartość 1. Jeżeli kolejka jest pusta, bit przyjmuje wartość 0 \cite{solomon2000threadScheduling}.

Dzięki takiej architekturze, problem wyboru następnego wątku sprowadza się do odnalezienia najstarszego ustawionego bitu (ang. \textit{most significant set bit}) w zmiennej \textit{Ready Summary}. Jest to operacja matematyczna, której złożoność czasowa wynosi $O(1)$. W architekturach sprzętowych z rodziny x86/amd64 poszukiwanie to realizowane jest za pomocą pojedynczej instrukcji procesora – BSR (\textit{Bit Scan Reverse}). Pozwala to w bardzo krótkim czasie zidentyfikować najwyższą niepustą kolejkę, pobrać z niej pierwszy oczekujący wątek i zainicjować przełączenie kontekstu \cite{reactosKeX}.


\section{Systemy z rodziny BSD na przykładzie FreeBSD}
W nowoczesnych wersjach systemu FreeBSD (począwszy od wydania 5.2) domyślnym podsystemem szeregowania zadań jest planista ULE \cite{roberson2003ule}. Zastąpił on historyczny algorytm wywodzący się z systemu 4.4BSD, który ze względu na globalną kolejkę uruchomieniową oraz konieczność okresowego przeliczania priorytetów wszystkich zadań w systemie (złożoność $O(n)$) nie skalował się efektywnie w środowiskach wieloprocesorowych (SMP) \cite{roberson2003ule, mckusick2014freebsd}. Nazwa ULE nie jest akronimem -- pochodzi od trzech ostatnich liter angielskiego słowa \textit{schedule} (powstałego po usunięciu znaku podkreślenia z nazwy pliku źródłowego \texttt{sched\_ule.c}) \cite{roberson2003ule, freebsdSchedUleC}. Analogicznie jak w systemie Windows, podstawową jednostką podlegającą szeregowaniu we FreeBSD nie jest cały proces (\texttt{struct proc}), lecz pojedynczy wątek jądra (\texttt{struct thread}), posiadający własną strukturę metadanych planisty \texttt{struct td\_sched} \cite{freebsdSchedUleC}.

Architektura FreeBSD definiuje 256 dyskretnych poziomów priorytetów (od 0 do 255). Zgodnie z klasyczną konwencją systemów UNIX, niższa wartość liczbowa oznacza wyższy priorytet wykonywania \cite{mckusick2014freebsd}. Przestrzeń ta podzielona jest na cztery główne klasy \cite{mckusick2014freebsd}:
\begin{itemize}
	\item \textbf{Poziomy 0--47 (Wątki przerwań i jądra):} Najwyższe priorytety nadawane podsystemom jądra, gwarantujące bezwzględne prawo wywłaszczania procesów przestrzeni użytkownika.
	\item \textbf{Poziomy 48--79 (Wątki czasu rzeczywistego):} Zadania o krytycznym znaczeniu czasowym (\textit{real-time}), których priorytety są nadawane bezpośrednio przez aplikacje i nie podlegają kwantowaniu czasu.
	\item \textbf{Poziomy 120--223 (Wątki podziału czasu):} Pasmo przeznaczone dla standardowych procesów użytkownika (\textit{timesharing}), w którym priorytet wyliczany jest dynamicznie przez planistę.
	\item \textbf{Poziomy 224--255 (Wątki bezczynności):} Zarezerwowane dla zadań wykonywanych wyłącznie przy braku jakichkolwiek innych gotowych wątków w systemie (\textit{idle threads}).
\end{itemize}

W momencie utworzenia nowego procesu za pomocą wywołania \texttt{fork()}, wątek potomny dziedziczy początkowy poziom priorytetu oraz skumulowaną historię interaktywności bezpośrednio po swoim procesie macierzystym (np. powłoce systemowej) \cite{roberson2003ule}. Sprawia to, że nowo powołane zadania zazwyczaj rozpoczynają pracę z wysokim priorytetem charakterystycznym dla aplikacji interaktywnych. W trakcie dalszego działania wątku w paśmie podziału czasu (120--223), jego priorytet ulega ciągłym zmianom w zależności od zachowania zadania oraz przypisanego mu przez użytkownika parametru \textit{nice} (z przedziału od -20 do +20) \cite{roberson2003ule, mckusick2014freebsd}.

Kluczowym mechanizmem modyfikującym priorytet w czasie pracy jest heurystyka oceny interaktywności, realizowana przez funkcję \texttt{sched\_interact\_score()} \cite{freebsdSchedUleC}. Planista monitoruje w strukturze \texttt{td\_sched} skumulowany czas dobrowolnego uśpienia wątku ($T_{\text{sleep}}$) oraz czas jego faktycznej pracy na procesorze ($T_{\text{run}}$) w obrębie ruchomego okna historii (ograniczonego stałą \texttt{SCHED\_SLP\_RUN\_MAX}) \cite{roberson2003ule, freebsdSchedUleC}. Oznaczając przez $M$ połowę maksymalnej skali interaktywności, wskaźnik $S$ wyznaczany jest zależnie od relacji obu czasów. Gdy czas snu przewyższa czas pracy ($T_{\text{sleep}} > T_{\text{run}}$), wskaźnik przyjmuje postać \cite{roberson2003ule}:

$$S = \left\lfloor \frac{M \cdot T_{\text{run}}}{T_{\text{sleep}}} \right\rfloor$$

Natomiast gdy dominuje czas obliczeń ($T_{\text{run}} \ge T_{\text{sleep}}$), wartość obliczana jest ze wzoru \cite{roberson2003ule}:

$$S = M + \left\lfloor \frac{M \cdot (T_{\text{run}} - T_{\text{sleep}})}{T_{\text{run}}} \right\rfloor$$

Im niższa jest wyliczona wartość $S$, tym wyższą interaktywnością cechuje się zadanie. Jeżeli wskaźnik $S$ spadnie poniżej zdefiniowanego progu, wątek zostaje uznany za interaktywny, a jego priorytet jest wyznaczany proporcjonalnie do wyniku interaktywności (z uwzględnieniem modyfikatora \textit{nice}) \cite{roberson2003ule}. Jeśli jednak proces (np. kompilator uruchomiony z interaktywnej powłoki) zacznie intensywnie obciążać jednostkę obliczeniową bez dobrowolnego zasypiania, wartość $T_{\text{run}}$ szybko zdominuje $T_{\text{sleep}}$. Przekroczenie progu interaktywności spowoduje degradację zadania do kategorii wsadowej (\textit{batch}). Dla takich wątków priorytet wewnątrz pasma podziału czasu rośnie (czyli staje się mniej korzystny) wraz z liczbą wyczerpywanych taktów zegara statystycznego (\texttt{statclock}) oraz dodatnią wartością \textit{nice} \cite{roberson2003ule, mckusick2014freebsd}.

\subsection{Wybór procesu i struktury danych}
Aby wyeliminować rywalizację o blokady między rdzeniami i osiągnąć złożoność czasową $O(1)$, scheduler ULE utrzymuje osobny zestaw kolejek uruchomieniowych \texttt{struct tdq} dla każdego procesora logicznego (w tablicy \texttt{tdq\_cpu}) \cite{freebsdSchedUleC}. Wewnątrz każdego kontenera (zestawu) \texttt{tdq} znajdują się trzy wyspecjalizowane struktury \texttt{struct runq}:
\begin{enumerate}
	\item \texttt{tdq\_realtime} -- przechowuje wątki przerwań, zadania czasu rzeczywistego oraz wątki podziału czasu zaklasyfikowane jako interaktywne \cite{mckusick2014freebsd, freebsdSchedUleC}.
	\item \texttt{tdq\_timeshare} -- gromadzi nieinteraktywne (wsadowe) wątki podziału czasu. Działa ona jako kolejka kalendarzowa (\textit{calendar queue}), wykorzystująca dwa rotujące indeksy: zapisu (\texttt{tdq\_idx}) oraz odczytu (\texttt{tdq\_ridx}). Wstawiane wątki umieszczane są w odległości od \texttt{tdq\_idx} zależnej od ich priorytetu, natomiast pobieranie odbywa się od pozycji \texttt{tdq\_ridx}, która przesuwa się z upływem czasu. Gwarantuje to, że każde zadanie wsadowe zostanie uruchomione podczas pełnego obrotu wskaźnika, co całkowicie zapobiega zjawisku głodzenia (\textit{starvation}) bez potrzeby globalnego przeliczania priorytetów \cite{mckusick2014freebsd, freebsdSchedUleC}.
	\item \texttt{tdq\_idle} -- gromadzi wątki bezczynności, uruchamiane wyłącznie wtedy, gdy dwie powyższe struktury są całkowicie puste \cite{freebsdSchedUleC}.
\end{enumerate}

Na najniższym poziomie każda z trzech instancji \texttt{struct runq} agreguje 256 poziomów priorytetów w 64 dwukierunkowe listy ogonowe (\texttt{rq\_queues[64]}). Samo odnalezienie niepustej kolejki wewnątrz wybranej struktury realizowane jest w czasie stałym $O(1)$ przez funkcję \texttt{runq\_choose()}, która odczytuje 64-bitowy wektor stanu \texttt{rq\_status}. Pozycja pierwszego ustawionego bitu bezpośrednio wskazuje indeks najważniejszej niepustej listy, z której pobierany jest pierwszy oczekujący wątek \cite{mckusick2014freebsd, freebsdRunqH}. Jest to bardzo podobne rozwiązanie optymalizacyjne jak w systemach z rodziny Windows.